% ID: SEL-D0019-OA-E01
% ESTADO: cotejada. Segunda pasada visual completa y borrador compilado revisados.
% Fuente: Selectividad/matematicas_aplicadas/reserva a/Examen-Matemáticas Aplicadas a las Ciencias Sociales.pdf
% SHA-256: 6fed0384ff6fcbaaa7098ab3f35ea5f881d1280502f050ada3a7f13af528392e
% Página: 1; opción A; ejercicio original 1.

\textbf{(2.5 puntos)} La capacidad máxima de trabajo de un taller que se dedica a la confección de pañuelos y corbatas es de 60 horas semanales. Cada pañuelo que confecciona le supone 2 horas de trabajo y le reporta un beneficio de 4 euros. En el caso de las corbatas son 3 horas y 6 euros respectivamente por unidad. Contrae el compromiso de que el número de corbatas confeccionadas más el doble del número de pañuelos debe ser, como mínimo, 28.

Con estas condiciones, \textquestiondown{}cuántas unidades de cada tipo de prenda debe confeccionar para obtener un beneficio económico máximo?
